\documentclass[a4paper]{report}

% --- FONT & NGÔN NGỮ ---
\usepackage{fontspec}
\setmainfont{Times New Roman}
\usepackage[vietnamese]{babel}
\usepackage{microtype} % Tự động tinh chỉnh khoảng cách để giảm lỗi tràn lề

% --- ÉP FONT 13pt TOÀN CỤC ---
\usepackage{fontsize}
\changefontsize[15.6]{13} % Cỡ chữ 13pt, base leading 15.6pt

% --- CANH LỀ CHUẨN A4 KHÓA LUẬN TỐT NGHIỆP TDTU ---
% Lề trái 3.5cm (đóng gáy sách), lề phải 2.0cm, lề trên 3.0cm, lề dưới 3.0cm
\usepackage[top=3.0cm, bottom=3.0cm, left=3.5cm, right=2.0cm, headheight=16pt, headsep=1.0cm, footskip=1.2cm]{geometry}

% --- DÃN DÒNG & THỤT ĐẦU DÒNG ĐOẠN VĂN ---
\usepackage{setspace}
\onehalfspacing % Dãn dòng 1.5 lines chuẩn
\usepackage{indentfirst}
\setlength{\parindent}{1.27cm} % Thụt đầu dòng 1.27cm (0.5 inch)
\setlength{\parskip}{0pt}      % Không dãn cách đoạn thừa thãi
\emergencystretch 3em          % Chống tuyệt đối lỗi tràn lề phải (overfull hbox)

% --- HEADER & FOOTER: ĐÁNH SỐ TRANG GIỮA TRÊN ---
\usepackage{fancyhdr}
\pagestyle{fancy}
\fancyhf{} 
\fancyhead[C]{\fontsize{10pt}{12pt}\selectfont\thepage} % Số trang căn giữa phía trên
\renewcommand{\headrulewidth}{0pt}
\renewcommand{\footrulewidth}{0pt}

% Đảm bảo trang đầu chương / trang plain vẫn có số trang căn giữa phía trên
\fancypagestyle{plain}{
    \fancyhf{} 
    \fancyhead[C]{\fontsize{10pt}{12pt}\selectfont\thepage} 
    \renewcommand{\headrulewidth}{0pt}
}

% --- TOÁN HỌC & HÌNH ẢNH ---
\usepackage{amsmath, amsfonts, amssymb}
\usepackage{graphicx}
\usepackage{tikz}
\usetikzlibrary{calc, positioning}
\usepackage{float}
\usepackage{longtable}
\usepackage{booktabs}
\usepackage{array}
\usepackage{multirow}
\usepackage{tabularx}
\usepackage{xltabular}
\newcolumntype{Y}{>{\centering\arraybackslash}X}
\usepackage{subcaption}
\graphicspath{{media/}}
\usepackage{caption}
\captionsetup{
    font={small},
    labelfont={bf},
    labelsep=colon,
    skip=6pt
}

\usepackage[hidelinks]{hyperref} % Liên kết mục lục sạch sẽ, không viền đỏ

% --- TÙY CHỈNH MỤC LỤC & DANH MỤC (TOC, LOF, LOT) ---
\usepackage{tocloft} 
\renewcommand{\cftdotsep}{0.5} % Dấu chấm mục lục san sát chuẩn Word

% Thêm chữ "CHƯƠNG " vào Mục lục
\renewcommand{\cftchappresnum}{CHƯƠNG } 
\renewcommand{\cftchapaftersnum}{. }    
\setlength{\cftchapnumwidth}{6.0em}     

% 1. Cấp Chương (Chapter): In đậm toàn hàng
\renewcommand{\cftchapfont}{\bfseries}
\renewcommand{\cftchapleader}{\bfseries\cftdotfill{\cftdotsep}}
\renewcommand{\cftchappagefont}{\bfseries}

% 2. Cấp Mục (Section - 1.1): Chữ thường
\renewcommand{\cftsecfont}{\normalfont}
\renewcommand{\cftsecleader}{\normalfont\cftdotfill{\cftdotsep}}
\renewcommand{\cftsecpagefont}{\normalfont}
\setlength{\cftsecnumwidth}{2.2em}

% 3. Cấp Tiểu mục (Subsection - 1.1.1): In nghiêng
\renewcommand{\cftsubsecfont}{\itshape}
\renewcommand{\cftsubsecleader}{\itshape\cftdotfill{\cftdotsep}}
\renewcommand{\cftsubsecpagefont}{\itshape}
\setlength{\cftsubsecnumwidth}{2.8em}

% Khoảng cách giữa các dòng trong Mục lục
\setlength{\cftbeforechapskip}{8pt} 
\setlength{\cftbeforesecskip}{5pt}
\setlength{\cftbeforesubsecskip}{4pt}

% Danh mục hình vẽ
\renewcommand{\cftfigpresnum}{Hình } 
\renewcommand{\cftfigaftersnum}{: }  
\setlength{\cftfignumwidth}{4.5em}   

% Danh mục bảng biểu
\renewcommand{\cfttabpresnum}{Bảng } 
\renewcommand{\cfttabaftersnum}{: }  
\setlength{\cfttabnumwidth}{4.5em}   

% Xóa khoảng trắng thừa giữa các chương trong danh mục hình/bảng
\usepackage{etoolbox}
\makeatletter
\AtBeginDocument{
  \patchcmd{\@chapter}{\addtocontents{lof}{\protect\addvspace{10\p@}}}{}{}{}
  \patchcmd{\@chapter}{\addtocontents{lot}{\protect\addvspace{10\p@}}}{}{}{}
}
\makeatother

% --- ĐỊNH DẠNG TIÊU ĐỀ CÁC CẤP (HEADING SPACING CHUẨN ĐẸP) ---
\usepackage{titlesec}
\setcounter{secnumdepth}{3} % Đánh số tới cấp 1.1.1.1

% Heading 1: Chương (\chapter) - 16pt Bold
\titleformat{\chapter}[hang]
  {\fontsize{16pt}{22pt}\selectfont\bfseries} 
  {CHƯƠNG \thechapter.} 
  {0.5cm} 
  {}
\titlespacing*{\chapter}{0pt}{0pt}{16pt} % Cách dưới 16pt

% Heading 2: Mục lớn (\section) - 14pt Bold, sát lề trái 0cm
\titleformat{\section}[hang]
  {\fontsize{14pt}{19pt}\selectfont\bfseries} 
  {\thesection}
  {0.4cm} 
  {}
\titlespacing*{\section}{0pt}{14pt}{6pt} % Sát lề trái

% Heading 3: Tiểu mục (\subsection - vd 3.1.1) - 13pt Bold Italic, thụt vào 0.635cm (nửa tab)
\titleformat{\subsection}[hang]
  {\fontsize{13pt}{18pt}\selectfont\bfseries\itshape} 
  {\thesubsection}
  {0.35cm} 
  {}
\titlespacing*{\subsection}{0.635cm}{10pt}{4pt} % Thụt vào bậc thang 0.635cm

% Heading 4: Tiểu mục con (\subsubsection - vd 3.1.1.1) - 13pt Italic, thụt vào 1.27cm (1 tab)
\titleformat{\subsubsection}[hang]
  {\fontsize{13pt}{18pt}\selectfont\itshape} 
  {\thesubsubsection}
  {0.3cm} 
  {}
\titlespacing*{\subsubsection}{1.27cm}{8pt}{3pt} % Thụt vào bậc thang 1.27cm

% --- KHOẢNG CÁCH CÔNG THỨC TOÁN HỌC ---
\AtBeginDocument{
  \setlength{\abovedisplayskip}{6pt}
  \setlength{\belowdisplayskip}{6pt}
  \setlength{\abovedisplayshortskip}{3pt} 
  \setlength{\belowdisplayshortskip}{3pt}
}

% --- TÀI LIỆU THAM KHẢO ---
\usepackage{csquotes} 
\DeclareQuoteAlias{english}{vietnamese}

\usepackage[backend=biber, style=ieee]{biblatex}
\addbibresource{backmatter/references.bib}
\usepackage{xurl} 
\let\parencite\cite 



\defbibheading{bibcenter}{
    \clearpage
    \phantomsection
    \addcontentsline{toc}{chapter}{TÀI LIỆU THAM KHẢO}
    \begin{center}
        {\fontsize{16pt}{22pt}\selectfont\bfseries TÀI LIỆU THAM KHẢO\par}
    \end{center}
    \vspace{12pt}
}

\defbibheading{subbib}{
    \noindent #1\par\vspace{6pt}
}

\defbibenvironment{bibliography}
  {\list
     {}
     {\setlength{\leftmargin}{0pt}
      \setlength{\itemindent}{1.27cm}
      \setlength{\itemsep}{0.5\baselineskip}
      \setlength{\parsep}{0pt}}}
  {\endlist}
  {\item}



% --- TÙY CHỈNH DANH SÁCH (ITEMIZE / ENUMERATE) ---
\usepackage{enumitem}
\setlist{
    topsep=3pt,
    partopsep=0pt,
    itemsep=2pt,
    parsep=0pt,
    labelindent=1.27cm,
    leftmargin=*,
    align=left,
    labelsep=0.4cm
}
\renewcommand{\labelitemi}{--}
\renewcommand{\labelitemii}{$-$}
% --- Lệnh đánh dấu chỗ cần điền số liệu thật, dùng thuần LaTeX (không cần gói ngoài) ---
% XOÁ lệnh \todo{...} này và thay bằng nội dung thật trước khi nộp báo cáo.
\newcommand{\todo}[1]{\textbf{[TODO: #1]}}